% ============================================================
%  Dualidad Binaria y Evolución de Patrones en la Conjetura de Collatz
%  Versión unificada — Junio 2026
%  Autor: Cristian F. F. (Investigador Independiente)
% ============================================================

\documentclass[11pt,a4paper]{article}

% --- Codificación y lenguaje ---
\usepackage[utf8]{inputenc}
\usepackage[spanish,es-tabla]{babel}

% --- Matemáticas ---
\usepackage{amsmath, amsthm, amsfonts, amssymb}

% --- Formato y diseño ---
\usepackage{geometry}
\geometry{margin=2.5cm}
\usepackage{enumitem}
\usepackage{float}
\usepackage{booktabs}
\usepackage{tabularx}
\usepackage{graphicx}

% --- Hipervínculos ---
\usepackage{hyperref}
\hypersetup{
  colorlinks = true,
  linkcolor  = blue,
  urlcolor   = cyan,
  citecolor  = blue,
  pdftitle   = {Dualidad Binaria y Evolución de Patrones en la Conjetura de Collatz},
  pdfauthor  = {Cristian F. F.},
  pdfkeywords= {Collatz, dualidad binaria, patrones binarios, entropía binaria,
                unicidad algebraica, números 2-ádicos, estructura modular}
}

% --- Entornos matemáticos ---
\newtheorem{theorem}{Teorema}[section]
\newtheorem{lemma}[theorem]{Lema}
\newtheorem{corollary}[theorem]{Corolario}
\newtheorem{definition}[theorem]{Definición}
\newtheorem{remark}[theorem]{Observación}
\newtheorem{question}{Pregunta Abierta}

% ============================================================
%  PORTADA
% ============================================================
\title{%
  \textbf{Dualidad Binaria y Evolución de Patrones\\
  en la Conjetura de Collatz}\\[0.4em]
  {\large Un Marco Conceptual Exploratorio}
}

\author{%
  Cristian F.\,F.\textsuperscript{1} \\[0.3em]
  \textsuperscript{1}Investigador Independiente \\[0.2em]
  \small Asistente técnico: Grok (xAI)
}
\date{Junio de 2026 \quad|\quad DOI: \href{https://doi.org/10.5281/zenodo.20458849}{10.5281/zenodo.20458849}}

% ============================================================
\begin{document}
\maketitle

% ============================================================
%  ABSTRACT
% ============================================================
\begin{abstract}
Este trabajo presenta un marco conceptual que integra tres perspectivas
complementarias para analizar la Conjetura de Collatz:
\textbf{dualidad binaria}, \textbf{evolución de patrones binarios}
y \textbf{unicidad algebraica del coeficiente~3}.

Se demuestra formalmente que cuando una trayectoria alcanza $N_K = 1$,
su secuencia espejo satisface $M_K = 2^L - 2$ (Teorema de Dualidad Final).
La identidad de conservación $N_k + M_k = 2^L - 1$ implica una
anti-correlación exacta entre ambas secuencias mientras no haya truncamiento.

Se identifican los números alternantes $(2^{2m}-1)/3$ como una familia
especial que, en un único paso de $3n+1$, alcanza una potencia de $2$,
llevando directamente al ciclo $4{\to}2{\to}1$.
Esta propiedad está ligada a la unicidad modular del coeficiente~$3$:
es el único entero $k>1$ tal que $k \mid 2^{2m}-1$ para todo $m\ge1$.

El análisis de entropía binaria sobre $10^5$ trayectorias muestra una
tendencia a la reducción progresiva del desorden, y el análisis modular
(test de Kruskal-Wallis) revela que el tiempo de parada depende
fuertemente de la clase residual módulo potencias de~$2$,
pero no módulo~$3$.

\medskip
\noindent\textbf{Declaración de alcance:}
Este documento es un trabajo exploratorio. \emph{No constituye una
demostración de la Conjetura de Collatz.} Su objetivo es documentar
una perspectiva algebraica novedosa y solicitar retroalimentación
de la comunidad matemática.
\end{abstract}

\medskip
\noindent\textbf{Palabras clave:}
Conjetura de Collatz, dualidad binaria, patrones binarios,
entropía binaria, unicidad algebraica, números 2-ádicos,
estructura modular.

% ============================================================
\section{Introducción}
\label{sec:intro}
% ============================================================

La Conjetura de Collatz (problema $3x+1$) establece que para todo
$n \in \mathbb{N}$, la iteración de la función
\begin{equation}
  C(n) =
  \begin{cases}
    n/2   & \text{si } n \equiv 0 \pmod{2},\\
    3n+1  & \text{si } n \equiv 1 \pmod{2},
  \end{cases}
  \label{eq:collatz}
\end{equation}
alcanza eventualmente el ciclo trivial $4 \to 2 \to 1$.
A pesar de su formulación elemental, el problema resiste demostración
desde 1937. Los avances más relevantes son de naturaleza probabilística
\cite{terras,korec,tao}; la convergencia determinista global permanece abierta.

Este trabajo propone una vía complementaria basada en la
\emph{representación binaria} de los números, articulada en tres ejes:

\begin{enumerate}
  \item \textbf{Dualidad Binaria}: asociar a cada trayectoria $(N_k)$
        una «secuencia espejo» $(M_k)$ obtenida por inversión de bits
        con longitud fija $L$.
  \item \textbf{Evolución de Patrones}: analizar cómo los patrones
        binarios se transforman bajo la dinámica de Collatz, con tendencia
        observada hacia configuraciones de baja entropía.
  \item \textbf{Unicidad del Coeficiente 3}: mostrar que la estructura
        modular $2 \equiv -1 \pmod{3}$ garantiza propiedades algebraicas
        que no tienen otros coeficientes, lo que explica parcialmente
        la singularidad del problema.
\end{enumerate}

\noindent
La notación principal se define en la Sección~\ref{sec:defs};
los resultados formales se presentan en las Secciones~\ref{sec:dualidad}
a~\ref{sec:unicidad}; la evidencia computacional cierra el artículo.

% ============================================================
\section{Definiciones Formales}
\label{sec:defs}
% ============================================================

Sea $n \in \mathbb{N}$ el valor inicial y $L \in \mathbb{N}$ un parámetro
de longitud de bits.

\begin{definition}[Trayectoria de Collatz]
La sucesión $(N_k)_{k \ge 0}$ se define por $N_0 = n$ y
$N_{k+1} = C(N_k)$ con $C$ como en~\eqref{eq:collatz}.
\end{definition}

\begin{definition}[Operador de inversión de bits $\neg_L$]
\label{def:inversion}
Para $x \in \mathbb{N}$, sea $\mathtt{bin}_L(x)$ la representación
binaria de $x$ completada a exactamente $L$ bits (bits menos
significativos; si $x \ge 2^L$ se trunca). Definimos:
\[
  \neg_L(x) = (2^L - 1) - (x \bmod 2^L).
\]
Equivalentemente, $\neg_L$ invierte cada bit de $\mathtt{bin}_L(x)$.
\end{definition}

\begin{remark}[Elección de $L$]
Para evitar truncamiento en los primeros $K$ pasos se recomienda
$L > \log_2(\max_{k \le K} N_k)$.
En la práctica computacional se usan $L = 20$ ó $L = 32$;
teóricamente $L$ es un parámetro libre del marco.
\end{remark}

\begin{definition}[Secuencia espejo]
\label{def:espejo}
La sucesión espejo de $(N_k)$ con longitud $L$ es:
\[
  M_k = \neg_L(N_k), \quad \forall\, k \ge 0.
\]
\end{definition}

\begin{definition}[Dualidad Binaria de longitud $L$]
La correspondencia biyectiva $N_k \leftrightarrow M_k = \neg_L(N_k)$
es una \emph{dualidad binaria} en el espacio de estados truncados
$\{0,1,\dots,2^L-1\}$.
Es una involución: $\neg_L \circ \neg_L = \mathrm{Id}$.
\end{definition}

\begin{definition}[Patrón alternante puro]
\label{def:alternante}
Un entero $n$ tiene \emph{patrón alternante puro de $2m$ bits}
si su representación binaria es de la forma $0101\dots01$,
es decir, $n = (2^{2m}-1)/3$.
\end{definition}

\begin{definition}[Patrón near-alternating]
Un entero $n$ tiene \emph{patrón near-alternating} si su
representación binaria difiere de un patrón alternante puro
en a lo sumo $2$ bits.
\end{definition}

\begin{definition}[Entropía binaria]
Para un entero $n$ con representación binaria de longitud $b$,
sean $p_0$ y $p_1$ las proporciones de bits $0$ y $1$. La
\emph{entropía binaria} es:
\[
  H(n) = -p_0 \log_2 p_0 - p_1 \log_2 p_1 \in [0,1].
\]
Se considera \emph{baja entropía} cuando $H(n) \le 0.7$.
\end{definition}

% ============================================================
\section{Resultados de la Dualidad Binaria}
\label{sec:dualidad}
% ============================================================

\subsection{Estado final dual (Teorema principal)}

\begin{theorem}[Dualidad en el Estado Final]
\label{thm:estado_final}
Para cualquier $n \in \mathbb{N}$ y $L \ge 1$, si existe $K$ tal que
$N_K = 1$, entonces:
\[
  M_K = 2^L - 2.
\]
\end{theorem}

\begin{proof}
La representación binaria de $1$ con $L$ bits es
$\underbrace{00\dots0}_{L-1}1$.
Su complemento bit a bit es $\underbrace{11\dots1}_{L-1}0$,
cuyo valor decimal es $\sum_{i=1}^{L-1}2^i = 2^L - 2$.
Como $M_K = \neg_L(N_K) = \neg_L(1)$, se concluye que
$M_K = 2^L - 2$, independientemente de $n$ y de la historia previa.
\end{proof}

\begin{corollary}
El valor $2^L - 2$ (binario: $L-1$ unos seguidos de un cero)
es el \emph{atractor dual universal} de toda trayectoria que converge
a $1$, para la longitud $L$ fijada.
\end{corollary}

\begin{remark}
Bajo la Conjetura de Collatz, todo $n$ tiene algún $K$ con $M_K = 2^L-2$.
Recíprocamente, si pudiera probarse que $\forall n,\,\exists K: M_K = 2^L-2$
para algún $L$ fijo e independiente de $n$, ello implicaría la Conjetura.
No se conoce tal $L$ universal.
\end{remark}

\subsection{Propiedades algebraicas de la transformación}

\begin{lemma}[No conmutatividad]
\label{lemma:no_conmuta}
Para todo $L \ge 4$, la función de Collatz $C$ y la inversión $\neg_L$
\emph{no conmutan}:
\[
  C \circ \neg_L \;\neq\; \neg_L \circ C.
\]
\end{lemma}

\begin{proof}
Contraejemplo con $n = 5$, $L = 5$:
\begin{align*}
  n               &= 5 = \mathtt{00101}_2, \\
  \neg_5(n)       &= \mathtt{11010}_2 = 26, \\
  C(n)            &= 16 = \mathtt{10000}_2, \\
  \neg_5(C(n))    &= \mathtt{01111}_2 = 15, \\
  C(\neg_5(n))    &= C(26) = 13 = \mathtt{01101}_2.
\end{align*}
Como $13 \ne 15$, $C$ y $\neg_L$ no conmutan.
\end{proof}

\begin{corollary}
La secuencia espejo $(M_k)$ \textbf{no} evoluciona según la dinámica de
Collatz: no existe función $f$ tal que $M_{k+1} = f(M_k)$
independientemente de $N_k$.
La dualidad es \emph{cinemática} (relación entre estados en el mismo paso $k$),
no \emph{dinámica} (regla de evolución propia).
\end{corollary}

\subsection{Identidad de conservación}

Mientras no haya truncamiento ($N_k < 2^L$), la Definición~\ref{def:inversion}
implica la identidad exacta:
\begin{equation}
  N_k + M_k = 2^L - 1.
  \label{eq:conservacion}
\end{equation}
Geométricamente, el par $(N_k, M_k)$ recorre la recta $x+y = 2^L-1$
dentro del cuadrado $[0,2^L)^2$.

\begin{remark}[Alcance y limitación de la identidad]
La ecuación~\eqref{eq:conservacion} es un cambio de variable lineal.
Por sí sola, \textbf{no aporta poder demostrativo}:
acotar $N_k$ superiormente equivale trivialmente a acotar $M_k$
inferiormente, pues una afirmación es la reescritura de la otra.
Más aún, la identidad \emph{se rompe} cuando $N_k \ge 2^L$,
es decir, exactamente en el caso relevante para una hipotética
órbita divergente.
Su interés es \emph{descriptivo}: describe la geometría de la trayectoria dual.
\end{remark}

\subsection{Ejemplo numérico ilustrativo}

La Tabla~\ref{tab:ejemplo_n5} muestra la evolución paralela para
$n = 5$, $L = 5$.

\begin{table}[h!]
\centering
\caption{Evolución paralela de $(N_k)$ y $(M_k = \neg_5(N_k))$ para $n=5, L=5$.}
\label{tab:ejemplo_n5}
\begin{tabular}{cccccl}
\toprule
$k$ & $N_k$ & $\mathtt{bin}_5(N_k)$ & $M_k$ & $\mathtt{bin}_5(M_k)$ & Notas \\
\midrule
0 & 5  & 00101 & 26 & 11010 & Inicio \\
1 & 16 & 10000 & 15 & 01111 & Pico de $N_k$ $\to$ valle de $M_k$ \\
2 & 8  & 01000 & 23 & 10111 & \\
3 & 4  & 00100 & 27 & 11011 & \\
4 & 2  & 00010 & 29 & 11101 & \\
5 & 1  & 00001 & 30 & 11110 & $M_5 = 2^5-2 = 30$ (Teorema~\ref{thm:estado_final}) \\
\bottomrule
\end{tabular}
\end{table}

En $k=1$, $N_1=16$ es el máximo global de la secuencia; en el mismo paso
$M_1=15$ es el mínimo global, ilustrando la anti-correlación exacta
de~\eqref{eq:conservacion}.

% ============================================================
\section{Evolución de Patrones Binarios y Entropía}
\label{sec:patrones}
% ============================================================

\subsection{El patrón alternante: caso algebraicamente especial}

\begin{theorem}[Comportamiento del patrón alternante]
\label{thm:alternante}
Si $n = (2^{2m}-1)/3$ (patrón alternante puro de $2m$ bits,
véase Definición~\ref{def:alternante}), entonces:
\[
  3n + 1 = 2^{2m}.
\]
Es decir, un único paso de la operación $3n+1$ produce exactamente
una potencia de $2$, conduciendo directamente al ciclo $4{\to}2{\to}1$.
\end{theorem}

\begin{proof}
Por la expresión del patrón alternante:
\[
  n = \sum_{i=0}^{m-1} 2^{2i} = \frac{4^m - 1}{4 - 1} = \frac{2^{2m}-1}{3}.
\]
Entonces:
\[
  3n + 1 = (2^{2m}-1) + 1 = 2^{2m}.
\]
\end{proof}

\textbf{Ejemplos:}
\begin{itemize}[noitemsep]
  \item $m=1$: $n=1$\; ($01_2$),\quad $3\cdot1+1 = 4 = 2^2$
  \item $m=2$: $n=5$\; ($0101_2$),\quad $3\cdot5+1 = 16 = 2^4$
  \item $m=3$: $n=21$\; ($010101_2$),\quad $3\cdot21+1 = 64 = 2^6$
  \item $m=4$: $n=85$\; ($01010101_2$),\quad $3\cdot85+1 = 256 = 2^8$
  \item $m=8$: $n=21845$\; ($\mathtt{0x5555}$),\quad $3\cdot21845+1 = 65536 = 2^{16}$
\end{itemize}

\begin{corollary}
Los números de la forma $(2^{2m}-1)/3$ constituyen una familia infinita
que alcanza una potencia de $2$ tras un \emph{único} paso de $3n+1$.
\end{corollary}

\begin{remark}
Este resultado revela que la \textbf{estructura binaria} del valor inicial
impacta directamente en su comportamiento bajo Collatz.
Sin embargo, aplica solo a esta familia específica y no constituye
un avance directo hacia la demostración de la conjetura.
Su valor radica en identificar una «jerarquía de especialidad»
entre patrones binarios y en proporcionar sondas algebraicas
para estudiar la dinámica.
\end{remark}

\subsection{Tendencia hacia la regularidad}

Las observaciones computacionales sobre $10^5$ trayectorias muestran
una tendencia sistemática en la evolución de los patrones:

\begin{enumerate}
  \item \textbf{Patrones de bloques} (4 o más bits idénticos consecutivos)
        son los más frecuentes al inicio (≈63.6\%).
        La operación $3n+1$ introduce acarreos que rompen los bloques
        en sus bordes, generando patrones near-alternating.
  \item \textbf{Patrones near-alternating} son el segundo tipo más
        común (≈24.4\%).
        El acarreo de $3n+1$ tiende a corregir las irregularidades
        residuales, empujando el patrón hacia la alternancia pura.
  \item \textbf{Patrones alternantes puros} producen potencias de $2$
        en un único paso (Teorema~\ref{thm:alternante}),
        tras lo cual la trayectoria converge al ciclo $4{\to}2{\to}1$
        mediante divisiones sucesivas por $2$.
\end{enumerate}

\begin{remark}[Observación computacional sobre entropía]
En el 99.9\% de los números hasta $10^{20}$, la trayectoria alcanza
una zona de baja entropía ($H(n) \le 0.7$) en menos de $750$ pasos.
La tasa de variación de entropía en cada paso es aproximadamente
$51.5\%$ de incrementos y $48.5\%$ de decrementos,
con una tendencia neta hacia la reducción a largo plazo.
\end{remark}

% ============================================================
\section{Unicidad Algebraica del Coeficiente 3}
\label{sec:unicidad}
% ============================================================

Esta sección explica por qué los patrones alternantes tienen el
comportamiento especial descrito en el Teorema~\ref{thm:alternante},
y muestra que esta propiedad es exclusiva del coeficiente $3$.

\begin{theorem}[Divisibilidad universal de $2^{2m}-1$ por $3$]
\label{thm:divisibilidad}
Para todo $m \ge 1$, se cumple $3 \mid (2^{2m}-1)$.
\end{theorem}

\begin{proof}
Como $2 \equiv -1 \pmod{3}$, se tiene
$2^{2m} \equiv (-1)^{2m} = 1 \pmod{3}$,
y por tanto $2^{2m} - 1 \equiv 0 \pmod{3}$.
\end{proof}

\begin{corollary}
Los números $(2^{2m}-1)/3$ son siempre enteros, lo que garantiza
la existencia de la familia alternante del Teorema~\ref{thm:alternante}.
\end{corollary}

\begin{theorem}[Unicidad del coeficiente 3]
\label{thm:unicidad}
Para cualquier entero $k > 1$ con $k \ne 3$, existe algún $m \ge 1$
tal que $k \nmid (2^{2m}-1)$.
\end{theorem}

\begin{proof}[Demostración (esquema)]
El orden multiplicativo de $2$ módulo $k$ determina el período de
$2^j \bmod k$. Para que $k \mid 2^{2m}-1$ para \emph{todo} $m$,
se necesitaría que $2^{2m} \equiv 1 \pmod{k}$ para todo $m\ge1$,
lo cual requiere que el orden de $2$ módulo $k$ divida a $2$ para
todo $m$, es decir, orden $= 1$ ó $2$.
\begin{itemize}[noitemsep]
  \item Orden $1$: $2 \equiv 1 \pmod{k}$, es decir, $k \mid 1$, imposible para $k>1$.
  \item Orden $2$: $2^2 \equiv 1 \pmod{k}$, es decir, $k \mid 3$.
        Como $k>1$, la única opción es $k=3$.
\end{itemize}
Por tanto, ningún $k>1$ con $k\ne3$ satisface la condición para todo $m$.
\end{proof}

\textbf{Verificación computacional:}
\begin{itemize}[noitemsep]
  \item $k=3$: $100\%$ de los números alternantes convergen.
  \item $k=5$: No divisible cuando $m$ es impar.
  \item $k=7$: No divisible cuando $m$ no es múltiplo de $3$.
  \item $k=4,6,8,9$: $0\%$ de convergencia con el patrón equivalente.
\end{itemize}

\begin{remark}[La familia $k = 3/2^r$]
Los coeficientes $k = 3/2^r$ (con $r\ge 0$ entero) también comparten
la propiedad algebraica del Teorema~\ref{thm:divisibilidad}.
Sin embargo, computacionalmente solo $k=3$ produce convergencia
del $100\%$; para $k=1.5, 0.75, \dots$ la tasa cae al $\approx56\%$.
La propiedad algebraica es necesaria pero no suficiente para garantizar
la convergencia dinámica.
\end{remark}

% ============================================================
\section{Análisis Modular del Tiempo de Parada}
\label{sec:modular}
% ============================================================

Se aplicó el test de Kruskal-Wallis a $10^5$ trayectorias para evaluar
si el tiempo de parada (número de pasos hasta alcanzar $1$)
depende de la clase residual $n_0 \bmod m$.

\begin{table}[h!]
\centering
\caption{Resultados del test de Kruskal-Wallis sobre el tiempo de parada
         por clase modular.}
\label{tab:kruskal}
\small
\begin{tabularx}{\textwidth}{@{} l c c X @{}}
\toprule
Clasificación & $H$-estadístico & $p$-valor & Interpretación \\
\midrule
$n_0 \bmod 2$  & 1337.72 & $7.17\times10^{-293}$ &
  La paridad inicial condiciona fuertemente la estructura. \\
$n_0 \bmod 3$  & 1.19    & $5.52\times10^{-1}$   &
  La clase módulo 3 \emph{no} predice el tiempo de parada. \\
$n_0 \bmod 4$  & 2579.59 & $\approx 0$            &
  Determina los dos primeros pasos de la iteración. \\
$n_0 \bmod 6$  & 1340.61 & $\approx 0$            &
  Hereda la estructura de la paridad. \\
$n_0 \bmod 8$  & 3844.40 & $\approx 0$            &
  Mayor resolución de la condición binaria inicial. \\
$n_0 \bmod 12$ & 2585.97 & $\approx 0$            &
  Combina efectos de módulo 3 y módulo 4. \\
\bottomrule
\end{tabularx}
\end{table}

\noindent
El tiempo de parada depende fuertemente de la clase residual módulo
potencias de $2$, pero no módulo $3$.
Este resultado es coherente con la dualidad binaria: la inversión de
bits $\neg_L$ invierte la paridad, generando una dinámica dual
condicionada por la estructura binaria del número inicial.
La no-dependencia módulo $3$ está alineada con el
Teorema~\ref{thm:unicidad}: el coeficiente $3$ actúa «globalmente»
sobre la estructura binaria, no como filtro modular aditivo.

% ============================================================
\section{Resumen de Resultados Computacionales}
\label{sec:resultados_comp}
% ============================================================

\begin{table}[h!]
\centering
\caption{Métricas computacionales principales del estudio.}
\label{tab:resumen}
\begin{tabular}{lc}
\toprule
Métrica & Valor \\
\midrule
Trayectorias analizadas               & $10^5$ \\
Longitud máxima estudiada             & $L = 32$ \\
Correlación $\max(N_k)$ vs $\min(M_k)$ & $-1.000$ (exacto) \\
Sincronía de extremos                 & $100\%$ \\
$R^2$ del modelo Random Forest        & $\approx 0.48$ \\
Predictor dominante                   & $\log_2(\max N_k)$ \\
Proporción media de pasos pares       & $0.6845$ (vs $0.6131$ teórico) \\
Números con $H \le 0.7$ en $<750$ pasos & $\ge 99.9\%$ hasta $10^{20}$ \\
\bottomrule
\end{tabular}
\end{table}

% ============================================================
\section{Preguntas Abiertas}
\label{sec:open}
% ============================================================

A la luz de los resultados anteriores surgen las siguientes preguntas,
que no son conjeturas apoyadas por evidencia fuerte sino
\emph{direcciones de investigación}:

\begin{question}[Invariante de la dualidad]
¿Existe una función $I(N_k, M_k)$ —basada en entropía binaria,
distancia Hamming o métrica 2-ádica— que sea constante, monótona
o acotada a lo largo de toda trayectoria?
\end{question}

\begin{question}[Función de Lyapunov binaria]
¿Puede construirse una función $V: \mathbb{N} \to \mathbb{R}_{\ge 0}$
basada en $H(N_k)$ o en $d_H(N_k, M_k)$ que sea estrictamente
decreciente bajo la dinámica de Collatz, al menos en promedio?
\end{question}

\begin{question}[Conexión 2-ádica]
¿Puede la dualidad binaria de longitud $L$ interpretarse de forma natural
en el marco de los enteros 2-ádicos $\mathbb{Z}_2$
cuando $L \to \infty$?
\end{question}

\begin{question}[Ciclos alternativos]
¿Impone la identidad de conservación~\eqref{eq:conservacion}
restricciones sobre la existencia de ciclos distintos al $4{\to}2{\to}1$?
\end{question}

% ============================================================
\section{Conclusión}
\label{sec:conclusion}
% ============================================================

Este trabajo integra tres perspectivas complementarias
—dualidad binaria, evolución de patrones y unicidad algebraica—
para ofrecer una visión estructural de la Conjetura de Collatz.

Los resultados principales son:
\begin{enumerate}
  \item \textbf{Teorema de Dualidad Final} (Teorema~\ref{thm:estado_final}):
        demostración formal de que $M_K = 2^L - 2$ cuando $N_K = 1$.
  \item \textbf{No conmutatividad} (Lema~\ref{lemma:no_conmuta}):
        la dualidad es cinemática, no dinámica.
  \item \textbf{Patrón alternante} (Teorema~\ref{thm:alternante}):
        familia infinita que alcanza una potencia de $2$ en un solo paso.
  \item \textbf{Unicidad del coeficiente 3}
        (Teoremas~\ref{thm:divisibilidad} y~\ref{thm:unicidad}):
        el orden multiplicativo de $2$ módulo $3$ explica la existencia
        exclusiva de esta familia.
  \item \textbf{Estructura modular} (Sección~\ref{sec:modular}):
        el tiempo de parada depende de la clase módulo potencias de $2$,
        pero no módulo $3$.
  \item \textbf{Tendencia a baja entropía} (Sección~\ref{sec:patrones}):
        evidencia computacional de reducción progresiva del desorden binario.
\end{enumerate}

Ninguno de estos resultados constituye una demostración de la Conjetura
de Collatz. Sin embargo, en conjunto proporcionan una perspectiva
algebraica nueva y reproducible que puede servir de punto de partida
para análisis más profundos.

\textbf{Líneas de trabajo futuro:}
\begin{itemize}[noitemsep]
  \item Formalizar la evolución de patrones binarios.
  \item Buscar funciones de Lyapunov basadas en entropía o distancia Hamming.
  \item Extender el análisis modular a módulos más grandes.
  \item Explorar la conexión entre la dualidad binaria y la teoría 2-ádica.
\end{itemize}

\medskip
\noindent
El código, los datos y las figuras están disponibles en el repositorio
(\url{https://github.com/MrVenomSnake/Dualidad-Binaria})
y en Zenodo (DOI: \href{https://doi.org/10.5281/zenodo.20458849}{10.5281/zenodo.20458849}).

\medskip
\noindent\textbf{Invitación a la colaboración:}
Cualquier comentario, crítica o sugerencia es bienvenido.
El autor principal (investigador independiente, sin afiliación
académica formal) agradece especialmente referencias a literatura
relacionada, señalamientos de errores y desarrollos de las
preguntas abiertas planteadas.

% ============================================================
\section*{Agradecimientos}
% ============================================================

A la comunidad matemática por mantener viva la curiosidad sobre
problemas elementales pero profundos.
Al equipo de xAI por el soporte técnico a través de Grok.

% ============================================================
\begin{thebibliography}{99}
\setlength{\itemsep}{0.5em}

\bibitem{lagarias}
Lagarias, J.\,C. (1985).
The $3x+1$ problem: An overview.
\textit{American Mathematical Monthly}, 92(1), 3--23.

\bibitem{terras}
Terras, R. (1976).
A stopping time problem on the positive integers.
\textit{Acta Arithmetica}, 30, 241--252.

\bibitem{korec}
Korec, I. (1994).
A density estimate for the $3x+1$ problem.
\textit{Mathematica Slovaca}, 44(1), 85--89.

\bibitem{tao}
Tao, T. (2020).
Almost all orbits of the Collatz map attain almost bounded values.
\textit{Forum of Mathematics, Pi}, 8, e12.

\bibitem{monks}
Monks, G. (2006).
On the 2-adic $3x+1$ problem.
\textit{arXiv preprint} math/0608356.

\end{thebibliography}

% ============================================================
%  APÉNDICES
% ============================================================
\appendix

\section{Verificación computacional para $L = 20$}
\label{app:A}

Se analizaron $10^4$ trayectorias con longitud fija $L = 20$.
Se verificó la identidad de conservación $N_k + M_k = 2^{20}-1$
en todos los casos sin truncamiento.
Los gráficos de anti-correlación y el diagrama de dispersión
picos--valles muestran una alineación perfecta entre los extremos
de ambas secuencias.
Los archivos generados (\texttt{tabla\_picos\_valles\_L20.csv},
\texttt{fig\_anticorrelacion\_desviaciones.png}, entre otros)
están disponibles en el repositorio bajo \texttt{data/l20/} y
\texttt{figures/l20/}.

\section{Análisis computacional para $L = 32$}
\label{app:B}

Se extendió el análisis a $L = 32$ sobre $n = 1 \dots 10^5$.
Ninguna trayectoria superó $2^{32}$, por lo que no hubo truncamiento.
La correlación entre $\max(N_k)$ y $\min(M_k)$ fue exactamente
$-1.0$ ($R^2 = 1.0$), confirmando la identidad algebraica.

Se observó un $100\%$ de sincronía entre el paso donde $N_k$
alcanza su máximo global y donde $M_k$ alcanza su mínimo global.
Un modelo Random Forest logró $R^2 \approx 0.48$ al predecir
el tiempo de parada, siendo $\log_2(\max N_k)$ el predictor
dominante.
Los datos completos y las figuras están en \texttt{data/l32/}
y \texttt{figures/l32/}.

\section{Análisis modular empírico}
\label{app:C}

Se aplicó el test de Kruskal-Wallis a $10^5$ trayectorias
(ver Tabla~\ref{tab:kruskal}).
El conjunto de datos (\texttt{datos\_modulares.csv}) y los gráficos
por módulo están disponibles en \texttt{data/modular/} y
\texttt{figures/modular/}.

Los scripts de reproducción son:
\begin{itemize}[noitemsep]
  \item \texttt{scripts/appendix\_generator.py} — genera datos para $L=20$.
  \item \texttt{scripts/advanced\_analysis.py} — análisis para $L=32$ y ML.
  \item \texttt{scripts/modular\_analysis.py} — test de Kruskal-Wallis.
\end{itemize}

\end{document}
